\section*{Chapter \ref{ch:g12:trade-and-trade-policy} --- Trade beyond Ricardo, and Trade Policy}
% figdata: figdata/grade-12/13-trade-and-trade-policy/out/ (tariffs, tariff_key, kestria_key); Kestria invented.

\begin{solution}{exo:g12:trade-and-trade-policy:1}
Free trade: $Q_d = 60$, $Q_s = 20$, imports 40. With the tariff, $P = 25$: $Q_d = 50$, $Q_s = 30$, imports 20.
\end{solution}

\begin{solution}{exo:g12:trade-and-trade-policy:2}
$50 \times 1.1 = 55$.
\end{solution}

\begin{solution}{exo:g12:trade-and-trade-policy:3}
32.2\,\% in 1992 and 2.2\,\% in 2022.
\end{solution}

\begin{solution}{exo:g12:trade-and-trade-policy:4}
Brazil, at 7.3\,\%: less than a seventh of India's 56.4\,\% in 1990.
\end{solution}

\begin{solution}{exo:g12:trade-and-trade-policy:5}
$(25 - 20) \times 20 = 100$.
\end{solution}

\begin{solution}{exo:g12:trade-and-trade-policy:6}
Revenue $5 \times 20 = 100$. \omterm{def:g10:what-prices-do:dwl}{Deadweight loss} $\tfrac12 \times 5 \times 10 + \tfrac12 \times 5 \times 10 = 50$.
\end{solution}

\begin{solution}{exo:g12:trade-and-trade-policy:7}
$P = 30$: $Q_d = Q_s = 40$, no imports and no revenue. \omterm{def:g10:what-prices-do:dwl}{Deadweight loss} $\tfrac12 \times 10 \times 20 + \tfrac12 \times 10 \times 20 = 200$: the tariff is prohibitive.
\end{solution}

\begin{solution}{exo:g12:trade-and-trade-policy:8}
The quota gives the same price, 25, the same quantities and the same \omterm{def:g10:what-prices-do:dwl}{deadweight loss}, 50. The rectangle of 100 goes to whoever holds the licences instead of the government.
\end{solution}

\begin{solution}{exo:g12:trade-and-trade-policy:9}
At 23: $Q_d = 54$, $Q_s = 26$, imports 28. Revenue $5 \times 28 = 140$, of which foreign sellers pay $2 \times 28 = 56$. \omterm{def:g10:what-prices-do:dwl}{Deadweight loss} $\tfrac12 \times 3 \times 6 \times 2 = 18$.
Net gain $56 - 18 = 38$, if the partner does not retaliate.
\end{solution}

\begin{solution}{exo:g12:trade-and-trade-policy:10}
Value added rises from $100 - 60 = 40$ to $110 - 60 = 50$: by 25\,\%, two and a half times the nominal tariff.
\end{solution}

\begin{solution}{exo:g12:trade-and-trade-policy:11}
Most of what buyers paid went to the American government as tariff revenue and to American producers as higher prices; only the two triangles, and the losses from
retaliation, are a net loss to the country.
\end{solution}

\begin{solution}{exo:g12:trade-and-trade-policy:12}
The prices of the targeted imports before duty did not fall: pass-through was complete, so American buyers paid the whole tariff and the United States gained nothing in its
\omterm{def:g10:ricardo-model:terms-of-trade}{terms of trade}.
\end{solution}

\begin{solution}{exo:g12:trade-and-trade-policy:13}
Cost: $10 \times (1 - 1.05^{-5})/0.05 \approx 43.3$. Savings: $8/0.05 = 160$ at the start of year 6, worth $160/1.05^5 \approx 125.4$ today. The savings are larger; but protection is the
right tool only if the savings would not come without it (lenders would not finance the learning), and if the protection really ends after five years.
\end{solution}

\begin{solution}{exo:g12:trade-and-trade-policy:14}
In the country rich in skilled labour, the scarce factor, unskilled labour, loses; in the country rich in unskilled labour, skilled labour, its scarce factor, loses.
\end{solution}

\begin{solution}{exo:g12:trade-and-trade-policy:15}
Domestic producers gain from any duty that raises their rivals' prices, while buyers' losses are spread thin; judging prices ``below cost'' leaves room for argument. Producers
therefore have reason to ask for duties often, and to spend on lobbying for them: \omterm{def:g11:public-choice-first-look:rent-seeking}{rent seeking}.
\end{solution}

\begin{solution}{pb:g12:trade-and-trade-policy:1}
\textbf{1.} Consumption 60, production 20, imports 40 million metres.

\textbf{2.} $40 \times 1.2 = 48$ kests.

\textbf{3.} Consumption 52, production 28, imports 24.

\textbf{4.} \omterm{def:g10:surplus-and-efficiency:consumer-surplus}{Consumer surplus} falls by $8 \times (60 + 52)/2 = 448$; \omterm{def:g10:surplus-and-efficiency:producer-surplus}{producer surplus} rises by $8 \times (20 + 28)/2 = 192$ million kests.

\textbf{5.} $8 \times 24 = 192$ million kests.

\textbf{6.} Production triangle $\tfrac12 \times 8 \times 8 = 32$; consumption triangle $\tfrac12 \times 8 \times 8 = 32$; \omterm{def:g10:what-prices-do:dwl}{deadweight loss} 64 million kests.

\textbf{7.} $448 = 192 + 192 + 64$.

\textbf{8.} 48 kests, the price at which Kestria's own buyers and producers leave a gap of 24.

\textbf{9.} $8 \times 24 = 192$ million kests: to the government if licences are auctioned, to the foreign exporters if they are given to them.

\textbf{10.} $64 + 192 = 256$ million kests.

\textbf{11.} Before: production 40, consumption 20, exports 20. After: production 34, consumption 26, exports 8.

\textbf{12.} $3 \times (40 + 34)/2 = 111$ million kests.

\textbf{13.} $3 \times (20 + 26)/2 = 69$ million kests.

\textbf{14.} $111 - 69 = 42$ million kests.

\textbf{15.} $64 + 42 = 106$ million kests.

\textbf{16.} Cloth producers and the treasury gain from the tariff, cloth buyers lose; after retaliation grain producers lose, grain buyers in Kestria gain a little.

\textbf{17.} Their gain of 192 is concentrated on a few firms, while the buyers' loss of 448 is spread over everyone.

\textbf{18.} Complete pass-through: the domestic price rose by the full tariff.

\textbf{19.} Part of the tariff would be paid by foreign sellers through a lower \omterm{def:g12:trade-and-trade-policy:world-price}{world price}, and Kestria might gain, unless its partners retaliated.

\textbf{20.} The \omterm{def:g10:what-prices-do:dwl}{deadweight loss} of the tariff is 64 million kests; after Brenland's retaliation, Kestria's net loss is 106 million kests.
\end{solution}
